% 04_preliminaries.tex
\section{Preliminaries}
\label{sec:prelim}

\textbf{Notation.} Let $\mathcal{V}$ denote the vocabulary of an
autoregressive language model, $V = |\mathcal{V}|$, and write
$\mathbf{t} = (t_1, \ldots, t_n) \in \mathcal{V}^n$ for a token
sequence of length $n$. The deployer and the verifier share a secret
key $\kappa \in \{0,1\}^*$ and the same tokenizer; the verifier
never sees the language model. Let
$\mathcal{H} : \{0,1\}^* \to \{0,1\}^{256}$ be SHA-256, modelled as
a random oracle in the security analysis (\autoref{thm:security}).
We write $[m] = \{0, \ldots, m-1\}$ and $z_\alpha$ for the one-sided
standard-normal quantile at level $\alpha$.

\textbf{Threat model.} The adversary observes watermarked text and
may paraphrase, translate, edit, splice, or delete tokens before
publication; we model these as a per-token modification rate
$\delta \in [0,1)$. The adversary is computationally bounded and
has no access to $\kappa$ or to the language model at detection
time. The verifier sees only token text and runs in $O(n)$ hash
operations.

\textbf{Detection problem and quality.} Detection is a binary
hypothesis test on text alone, $H_0$ (non-watermarked) vs.\ $H_1$
(watermarked under key $\kappa$), at a fixed false-positive rate
$\alpha$, since false positives are operationally the costly error.
The watermark distribution must remain fluent under standard greedy
or low-temperature decoding, which we control via the per-token KL
divergence to the original next-token distribution and report as
perplexity inflation (\autoref{sec:method}).
